Abbreviate \(M_n=2^{dJ_z}\), \(K=K_{J_y,J_z,n}\), \(A=A_{J_y,J_z,n}\), \(V=V_{L,n}\), and \(S=S_n=A+\sqrt{nV}\).  The uniform second-order expansion and the corrected bias premise imply
\[
n\{\widehat\theta^{\mathrm{an}}_{J_y,J_z,n}-\theta_n-B^y_{J_y,n}-B^z_{J_z,n}\}
=n\mathbb L_{J_y,J_z,n}+n\mathbb U_{J_y,J_z,n}+o_P(S).
\tag{24}
\]
Indeed, \(n\rho_n/\sqrt n\asymp\sqrt{nV}\le S\) by the variance assertion in \(\mathrm{thm{:}uniform\text{-}second\text{-}order\text{-}expansion}\), \(A\le S\), and the two biases multiplied by \(n\) are \(O(S)\) by the displayed premise.

We next identify the local linear term, including the split constants.  Let \(s_{a,n}\) be the exact centered first derivative of the conditional transport functional and let \(s^{(2)}_{a,J_y,J_z,n}\) be its projected inverse-elliptic linearization.  Define the conditional distribution functions \(F_{a,n}(y\mid z)=\int_0^yq_{a,n}(u,z)du\), their generalized inverses \(F_{a,n}^{-1}(\cdot\mid z)\), the optimal map \(T_{n,z}=F_{1,n}^{-1}(\cdot\mid z)\circ F_{0,n}(\cdot\mid z)\), and \(H_n(y,z)=\int_0^y\Delta_n(u,z)du\).  Expanding the identity \(F_{1,n}\{T_{n,z}(y)\mid z\}=F_{0,n}(y\mid z)\) by the mean-value theorem, and using \(c\le q_{a,n}\le C\), gives
\[
T_{n,z}(y)-y
=-\frac{H_n(y,z)}{\bar q_n(y,z)}
+O\{\|\Delta_n\|_\infty |H_n(y,z)|\}.
\tag{25}
\]
The inverse-elliptic energy defining \(\rho_n\) is equivalent, by the same density bounds, to \(\|H_n\|_{L^2(dy\,dR_n)}^2\).  The outcome antiderivative gains one wavelet order, while the covariate projection retains the declared \(s_z\) smoothness.  The standard coefficient estimates following from \(\mathrm{ass{:}anisotropic\text{-}conditional\text{-}smoothness}\) and \(\mathrm{ass{:}wavelet\text{-}regularity}\), applied separately in the two tensor coordinates, are therefore
\[
\|(I-P_{J_y})H_n\|_2\lesssim2^{-(s_y+1)J_y},
\qquad
\|(I-P_{J_z})H_n\|_2\lesssim2^{-s_zJ_z}.
\tag{26}
\]
Equations (25)--(26) prove
\[
\max_a\|s_{a,n}-s^{(2)}_{a,J_y,J_z,n}\|_{L^2(P_{a,n})}
\le C\{\rho_n\|\Delta_n\|_\infty+2^{-(s_y+1)J_y}+2^{-s_zJ_z}\}.
\tag{27}
\]

For fold \(e\), put
\[
\bar\Phi_{a,e}=m^{-1}\sum_{i\in I_{a,e}}\Phi_{J_y,J_z}(W_{a,i,n}),
\quad
Z_{e,n}=\Gamma^{-1/2}\sqrt m\{\bar\Phi_{1,e}-\bar\Phi_{0,e}-\delta\},
\tag{28}
\]
where \(\Gamma=\Gamma_{J_y,J_z,n}\) and \(\delta=\delta_{J_y,J_z,n}\).  The two vectors in (28) are independent, centered, and have covariance identity by \(\mathrm{ass{:}two\text{-}arm\text{-}sampling}\).  Define
\[
G_n=2\sum_{e=1}^2\sqrt m\,\delta^{\top}B\Gamma^{1/2}Z_{e,n}.
\tag{29}
\]
If \(Z_{e,k}=v_k^{\top}Z_{e,n}\), diagonalization of \(K=\Gamma^{1/2}B\Gamma^{1/2}\), together with the definition of \(\eta_{k,n}\), gives the exact identities
\[
G_n=2\sum_{e=1}^2\sum_k\nu_k\eta_{k,n}Z_{e,k},
\qquad
\operatorname{Var}(G_n)=8\sum_k\nu_k^2\eta_{k,n}^2.
\tag{30}
\]
An empirical average of a centered score error of \(L^2\)-norm \(r_n\), multiplied by \(n\), has \(L^2\)-norm at most \(C\sqrt n r_n\).  Thus (27) yields
\[
\|n\mathbb L-G_n\|_2
\le C\sqrt n\{\rho_n\|\Delta_n\|_\infty+2^{-(s_y+1)J_y}+2^{-s_zJ_z}\}=o(S).
\tag{31}
\]
For the first term, use \(\rho_n\sqrt n\asymp\sqrt{nV}\le S\) and \(\|\Delta_n\|_\infty\to0\).  For \(x_n=2^{-J_y}\), the corrected premise gives
\[
\frac{n x_n^{2s_y+2}}{S^2}
=\frac{n x_n^{2s_y+1}}S\frac{x_n}S=o(1),
\]
and similarly \(n2^{-2s_zJ_z}/S^2=O(S^{-1})=o(1)\).  Here \(S\ge A\asymp M_n^{1/2}\to\infty\) by \(\mathrm{prop{:}anisotropic\text{-}spectrum\text{-}bias}\).  Equations (30)--(31) and \(n\operatorname{Var}(\mathbb L)=V\) imply
\[
8\sum_k\nu_k^2\eta_{k,n}^2=nV+o(S^2).
\tag{32}
\]

Put \(Q_n=n\mathbb U_{J_y,J_z,n}\).  The proved \(\mathrm{lem{:}growing\text{-}kernel\text{-}contractions}\) applies: \(M_n\le D_{J_y,J_z}=o(n^2)\), and the one-dimensional Green bound used in its proof makes the projected score summands in (29) uniformly bounded.  Its martingale differences sum to \(G_n+Q_n\), obey Lindeberg, and have predictable conditional variance converging in \(L^2/S^2\) to their unconditional variance.  Hoeffding orthogonality, (14), and (32) give
\[
E(G_n+Q_n)^2=A^2+nV+o(S^2).
\tag{33}
\]
For fixed \(t\), conditionally use
\[
|e^{ix}-1-ix+x^2/2|\le Cx^2\min\{|x|,1\}.
\]
The Lindeberg condition makes the sum of the conditional Taylor remainders vanish.  The predictable-variance conclusion and the Lipschitz continuity of the exponential then replace the random quadratic term by (33).  Iterating the conditional expectations over the ordered observations proves
\[
E\exp\{it(G_n+Q_n)/S\}
-\exp\!\left[-\frac{t^2}{2S^2}\{A^2+nV+o(S^2)\}\right]\longrightarrow0.
\tag{34}
\]

Define the finite Gaussian-quadratic proxy
\[
\mathcal Q_n^G=\sum_{e=1}^2\sum_k\nu_k
\{(\xi_{e,k}+\eta_{k,n})^2-(1+\eta_{k,n}^2)\}.
\tag{35}
\]
For \(c_{e,k}=\nu_k/S\) and \(l_{e,k}=2\nu_k\eta_{k,n}/S\), direct one-dimensional Gaussian integration gives
\[
\log E e^{it\mathcal Q_n^G/S}
=\sum_{e,k}\left[-itc_{e,k}-\frac12\log(1-2itc_{e,k})
-\frac{t^2l_{e,k}^2}{2(1-2itc_{e,k})}\right].
\tag{36}
\]
By (32) and the spectral proposition,
\[
\sum_{e,k}l_{e,k}^2=\frac{nV+o(S^2)}{S^2},
\qquad
\max_k\frac{|\nu_k|}{S}\le\frac{\max_k|\nu_k|}{A}lesssim M_n^{-1/2}\to0.
\tag{37}
\]
Expanding (36), the centered-quadratic remainder is bounded by
\(C\max_k|\nu_k|S^{-1}\sum_{e,k}\nu_k^2S^{-2}=o(1)\), and the noncentral linear remainder by
\(C\max_k|\nu_k|S^{-1}\sum_{e,k}l_{e,k}^2=o(1)\).  Since \(A^2=4\sum_k\nu_k^2\), this proves
\[
\log E e^{it\mathcal Q_n^G/S}
=-\frac{t^2}{2S^2}(A^2+nV)+o(1).
\tag{38}
\]

The variance \(\sigma_n^2=(A^2+nV)/S^2\) lies in \([1/2,1]\).  Every subsequence has a further subsequence on which \(\sigma_n^2\) converges.  Equations (34) and (38) then make both \((G_n+Q_n)/S\) and \(\mathcal Q_n^G/S\) converge on that further subsequence to the same centered normal law.  Their second moments are uniformly bounded by (33) and (38), hence the subsequence criterion for bounded-Lipschitz distance yields
\[
d_{\mathrm{BL}}\{\mathcal L((G_n+Q_n)/S),\mathcal L(\mathcal Q_n^G/S)\}\to0
\]
and convergence of each law to \(\mathcal N(0,\sigma_n^2)\) in bounded-Lipschitz distance.  Equations (24) and (31) replace \(G_n+Q_n\) by the normalized estimator, proving both displayed conclusions of the proposed theorem.

At conditional equality, \(\Delta_n=0\), hence \(\delta=0\), every \(\eta_{k,n}=0\), \(V=0\), and \(\mathbb L=0\).  Taking \(S=A\) in (36)--(38) gives
\[
n\mathbb U_{J_y,J_z,n}/A\Rightarrow\mathcal N(0,1).
\]
The decisive condition is the proved ratio \(\max_k|\nu_k|/A\lesssim M_n^{-1/2}\to0\), so no fixed or subsequential non-Gaussian chaos survives.

Finally consider a fixed unequal law.  Its exact centered Kantorovich scores are bounded by compact support and the positive density lower bound.  Applying the same conditional characteristic-function expansion as above to their independent sum gives
\(\mathbb L/\sqrt{V/n}\Rightarrow\mathcal N(0,1)\).  The assumption \(\sqrt{nV}/A\to\infty\), together with (14), gives
\[
\mathbb U=o_P\{\sqrt{V/n}\}.
\]
The assumed corrected bias and remainder bounds remove the other expansion terms, and \(\widehat V_{L,n}/V_{L,n}\to_P1\) permits studentization.  This proves the last displayed central limit theorem.